\documentclass[11pt]{article}
\usepackage[utf8]{inputenc}
\usepackage{graphicx}
\usepackage[margin=2cm]{geometry}      % default margins are around 3.5–4 cm
\usepackage[style=ieee, backend=biber, maxbibnames=1]{biblatex}
\usepackage[htt]{hyphenat}
\usepackage{enumitem}
\setlist{nosep}                        % tighter itemize/enumerate

\usepackage[compact]{titlesec}         % less space around section headings

\renewcommand*{\bibfont}{\small}  
\usepackage{booktabs}
\usepackage{tabularx}
\usepackage{float}


\title{Term Rewriting Systems}
\author{Allan Ken Miyazono Ushijima \and Luis Fregoso}

\date{}
\begin{document}

\maketitle
\section{Rewrite Rules}
\begin{itemize}
    \item \textbf{binaryOperatorp(): }It is a rule that recognizes the operators as facts of binaryOperatorp.
    \item \textbf{unaryOperatorp(): }It is a rule that recognizes $\neg$ as a fact of unaryOperatorp.
    \item \textbf{constantp(P): }It is a rule that recognizes a propositional constant.
    \item \textbf{propositionp(P): }It is a rule that recognize a proposition. It has to either be a constantp, a proposition with a $\neg$, or two propositions connected with a binaryOperatorp.
    \item \textbf{formula2cnf(P, \_): }It is the a rule that check if it is appropiate language (propositionp(P)), transforms it to conjunctive normal form (cnf()), splits propositions into conjuncts ($\land$), and splits conjuncts by the $\lor$ (maplist()).
    \item \textbf{cnf(P, CNF): }It is the rule where the main transformation occurs: \textit{elim(P, P1)} that removes $\leftrightarrow$ and $\rightarrow$, \textit{nnf(P1, P2)} that pushes down the $\neg$ down to atoms, and \textit{dist} that deals with the distributive laws.
    \item \textbf{elim(P, P1): }It is the rule that eliminates the double implication ($\leftrightarrow$) and implication ($\rightarrow$). The biconditional doesn't transform to implication, but to conjuctions and disjunctions directly.
   
    $ P \leftrightarrow Q \equiv (P \rightarrow Q) \land (Q \rightarrow P)$
    
    $ P \leftrightarrow Q \equiv (\neg P \lor Q) \land (\neg Q \lor P) $
    
    $P \rightarrow Q \equiv \neg P \lor Q$

    \item \textbf{nnf(P, P1): }This deals with either the double negation, or De Morgan's laws: 
    
    $\neg (\neg P) \equiv P$

    $\neg(P \lor Q) \equiv \neg P \land \neg Q $ 
    
    $ \neg(P \land Q) \equiv \neg P \lor \neg Q$

    \item \textbf{dist(P, CNF): }This rule formats the proposition to CNF using distribution laws. 
    
    $P \lor (Q \land R) \equiv (P \lor Q) \land (P \lor R)$

    $(Q \land R) \lor P \equiv (P \lor Q) \land (P \lor R)$

    \item \textbf{}
\end{itemize}

\section{Control Strategies}

\section{Correctness Argument}

\section{Testing Evidence}

\section{Known Limitations}

\section{Contributions}

\end{document}
